%%fssp_ncor_specacr
L'algorithme et les formules \'etant complexes
seul le principe de calcul est pr\'esent\'e ici.\\
Selon Brenguier-89 \cite{Brenguier-89}, le
temps de transit moyen d'une gouttelettes   
s'exprime comme une fonction~: \\ 
\begin{equation}
\overline{T}_c = f_{\overline{T}} ( N_c, \tau, \tau_1, \tau_2, \tau_r )
\label{eqn:tbar}  
\end{equation}
o\`u~:\\
$\overline{T_c} $ : Temps de transit moyen \\
$N_c$ : nombre de gouttes corrig\'e des
co\"{\i}ncidences \\
$\tau $ : temps de transit d'une particule \\
\\
On pose $\tau = \tau_{spec}$, $\overline{T}_c = \overline{T}$.
L'\'equation (\ref{eqn:tbar}) ne poss\`ede plus qu'une inconnue $N_c$. 
Elle est r\'esolue par it\'erations successives et donne  
$N_{loc}$.  $H_c$ est calcul\'e par
$H_c = A_m/A$ o\`u 
$A = f_A(N_{loc}, \tau_{spec}, \tau_1, \tau_2, \tau_r )$.\\ 
Si $Reset\times f_e \le Crt_0$ alors aucune correction
est calcul\'ee et $N_{loc}=0$ et $H_c=0$.
